\corrigesChapitre{8}{Nombres premiers et petit théorème de Fermat}
\begin{corrige}{1}{Premiers ou composés}
\begin{questions}
\item $1$ n’est ni premier ni composé ; $2,29$ sont premiers ; $21$ est composé.
\item $\sqrt{29}<6$ : tout diviseur premier d’un composé serait parmi $2,3,5$. Aucun ne divise $29$.
\item $84=2^2\times3\times7$.
\end{questions}
\end{corrige}
\begin{corrige}{2}{Décomposition et diviseurs}
\begin{questions}
\item $(3+1)(2+1)(1+1)=24$, en choisissant indépendamment chaque exposant.
\item Les exposants minimaux donnent $2^2\times3=12$.
\item Les maximaux donnent $2^3\times3^2\times5\times7=2520$.
\end{questions}
\end{corrige}
\begin{corrige}{3}{Fermat direct}
\begin{questions}
\item $1$, car $13$ est premier et ne divise pas $2$.
\item $100=12\times8+4$, donc $2^{100}\equiv2^4\equiv3$.
\item L’hypothèse $13\nmid a$ échoue ; $13^{12}\equiv0$, pas $1$.
\end{questions}
\end{corrige}
\begin{corrige}{4}{Un chiffrement}
\begin{questions}
\item A donne $8$, soit I ; B donne $13$, soit N.
\item $5\times21=105\equiv1\pmod{26}$.
\item $x\equiv21(y-8)\pmod{26}$.
\end{questions}
\end{corrige}
\begin{corrige}{5}{Primalité de $223$}
\begin{questions}
\item $\sqrt{223}\simeq14{,}93$.
\item Tester $2,3,5,7,11,13$. Les restes sont respectivement $1,1,3,6,3,2$.
\item Aucun ne divise $223$, donc $223$ est premier.
\end{questions}
\end{corrige}
\begin{corrige}{6}{Crible}
\begin{questions}
\item Les multiples $kp$ avec $2\leq k<p$ ont déjà un facteur premier plus petit que $p$ et ont été éliminés.
\item $2,3,5,7,11,13,17,19$.
\item Dans les deux cas, la liste rendue est vide ; zéro et un sont marqués faux.
\end{questions}
\end{corrige}
\begin{corrige}{7}{Décomposition algorithmique}
\begin{questions}
\item La liste est $[2,2,3,3,5]$.
\item $2^2\times3^2\times5=180$.
\item Sinon il posséderait un diviseur au plus sa racine, déjà testé avant l’arrêt ; contradiction.
\end{questions}
\end{corrige}
\begin{corrige}{8}{Un inverse par Fermat}
\begin{questions}
\item $a^{p-2}$ convient car $a\times a^{p-2}=a^{p-1}\equiv1$.
\item $3^5\equiv5\pmod7$, donc inverse $5$.
\item Bézout donne aussi $3\times5-7\times2=1$. Fermat suppose le module premier ; Bézout fonctionne pour tout module premier avec $a$.
\end{questions}
\end{corrige}
\begin{corrige}{9}{Un faux test de primalité}
\begin{questions}
\item $2^{10}=1024=3\times341+1$.
\item $2^{340}=(2^{10})^{34}\equiv1$.
\item On ne doit pas conclure que $341$ est premier. Une congruence de Fermat vérifiée est une condition nécessaire dans ce cadre, pas suffisante.
\end{questions}
\end{corrige}
\begin{corrige}{10}{L’infinité, sans faux raccourci}
\begin{questions}
\item $N=30031$.
\item $30031=59\times509$.
\item La preuve exige un diviseur premier de $N$, pas que $N$ lui-même soit premier. Aucun premier de la liste initiale ne divise $N$.
\end{questions}
\end{corrige}
\begin{corrige}{11}{RSA miniature, guidé}
\begin{questions}
\item $7^3=343\equiv13\pmod{55}$.
\item $81=2\times40+1$ ; l’exposant de déchiffrement proposé est $27$.
\item Le calcul rend $7$, le message initial. Ce petit exemple illustre le mécanisme ; ses nombres sont trop petits pour une sécurité réelle.
\end{questions}
\end{corrige}
\begin{corrige}{12}{Racine rationnelle, prolongement}
\begin{questions}
\item $a_np^n+a_{n-1}p^{n-1}q+\cdots+a_0q^n=0$.
\item Les autres termes sont divisibles par $q$, donc $q\mid a_np^n$ ; comme $q$ est premier à $p^n$, Gauss donne $q\mid a_n$.
\item Isoler le dernier terme : $p\mid a_0q^n$ ; $p$ et $q^n$ sont premiers entre eux, donc Gauss donne $p\mid a_0$.
\end{questions}
\end{corrige}
